% vim: set ft=tex tabstop=4 shiftwidth=4 noexpandtab:

\input{../../include/latex/tikz/header-light.tex}

% opening %{{{1

\begin{document}
\begin{tikzpicture}[scale=1.0]

% parameters %{{{1

	\def\len{0}
	\def\gth{8}

	\def\ang{45}

	\def\ratio{0.0}
	\def\fraction{1.5}

	%\tikzmath{
		%\var = 2.1;
	%}

% coordinates %{{{1

	\coordinate (DX) at (180:\len);
	\coordinate (DY) at (0:\gth);

	\coordinate (FX) at ({\ang + 180}:\len);
	\coordinate (FY) at (\ang:\gth);

	\coordinate (S) at (0,0);
	\coordinate (J) at (0:{0.4*\gth});
	\coordinate (K) at (\ang:{0.4*\gth});

	\coordinate (A) at ({\ang/2}:{0.7*\gth});

	\coordinate (SA) at ($ (A) - (S) $);
	\coordinate (AX) at ($ (S) - \ratio*(SA) $);
	\coordinate (AY) at ($ (S) + \fraction*(SA) $);

	\coordinate (O) at ($(DX)!(A)!(DY)$);
	\coordinate (P) at ($(FX)!(A)!(FY)$);

% intersections %{{{1

	%\begin{pgfinterruptboundingbox}
		%\path[name path=line1] (A) -- (C);
		%\path[name path=line2] (B) -- (D);
		%\path[name intersections={of=line1 and line2, by=S}];
	%\end{pgfinterruptboundingbox}

	%\begin{pgfinterruptboundingbox}
		%\path[name path=line] (A) -- (B);
		%\path[name path=circle] (C) circle (\radius);
		%\path[name intersections={of=line and circle, by={S,J}}];
	%\end{pgfinterruptboundingbox}

	%\begin{pgfinterruptboundingbox}
		%\path[name path=circleA] (A) circle (\radius);
		%\path[name path=circleB] (B) circle (\radius);
		%\path[name intersections={of=circleA and circleB, by={S,J}}];
	%\end{pgfinterruptboundingbox}

% distance %{{{1

	%\path let \p1 = ($ (B) - (A) $)
	%in \pgfextra{
		%\pgfmathsetmacro{\distAB}{veclen(\x1,\y1)/1cm}
		%\global\let\distAB\distAB
	%};

% points, dots, vertices %{{{1

	\fill (S) circle (0.4mm);
	%\fill (J) circle (0.4mm);
	%\fill (K) circle (0.4mm);

	\fill (A) circle (0.4mm);

	\fill (O) circle (0.4mm);
	\fill (P) circle (0.4mm);

	%\fill (DX) circle (0.4mm);
	%\fill (DY) circle (0.4mm);

	%\fill (FX) circle (0.4mm);
	%\fill (FY) circle (0.4mm);

% segments, sides, lines %{{{1

	\draw (DX) -- (DY);
	\draw (FX) -- (FY);

	\draw (AX) -- (AY);

	\draw (A) -- (O);
	\draw (A) -- (P);

	%\draw (A) -- (B) -- (C) -- cycle;

% circles %{{{1

	%\draw (O) circle (\radius);

	%\path[draw] let \p1 = ($ (A) - (O) $)
	%in (O) circle ({veclen(\x1,\y1)});

% circular arcs %{{{2

	%\draw (A) ++(\ang:\radius) arc (\ang:\gle:\radius);

% points, dots, vertices labels %{{{1

	\node[below left] at (S) {$S$};
	%\node[below] at (J) {$J$};
	%\node[above left] at (K) {$K$};
	\node[above] at (A) {$A$};
	\node[below] at (O) {$O$};
	\node[above left] at (P) {$P$};

% segments, sides, lines labels %{{{1

	\node[above] at ($ (DX)!0.9!(DY) $) {$d$};
	\node[above left] at ($ (FX)!0.9!(FY) $) {$f$};
	%\node[above left] at ($ (AX)!0.9!(AY) $) {$b$};

% segments, sides, lines marks %{{{1

	%\tkzMarkSegments[mark=||, size=2pt](A,O)
	%\tkzMarkSegments[mark=||, size=2pt](A,P)

% circles labels %{{{1

	%\node at ($ (O) + (45:\radius + 0.3) $) {$\mathscr{C}$};

% circular arcs labels %{{{2

	%\node at ($ (O) + (45:\radius + 0.3) $) {$\mathscr{C}$};

% angles labels %{{{1

	\pic[draw, ->, "$\alpha$", angle radius=1.9cm, angle eccentricity=1.2]
	{angle = O--S--A};

	\pic[draw, ->, "$\beta$", angle radius=1.8cm, angle eccentricity=1.2]
	{angle = A--S--P};

% right angles markers %{{{2

	\pic [draw, angle radius=7pt, angle eccentricity=1]
	{right angle = S--O--A};

	\pic [draw, angle radius=7pt, angle eccentricity=1]
	{right angle = S--P--A};

% closing %{{{1

\end{tikzpicture}
\end{document}
